\chapter*{Hướng Dẫn Sử Dụng Sách}
\addcontentsline{toc}{chapter}{Hướng Dẫn Sử Dụng Sách}
\markboth{Hướng Dẫn Sử Dụng Sách}{Hướng Dẫn Sử Dụng Sách}

\begin{truyen}{Đọc Cuốn Sách Này Như Thế Nào?}{How to Read This Book}
\chuhoa{M}{ỗi chương trong sách có ba lớp. Lớp đầu là một câu chuyện gần với đời sống của một người thầy, người cha, người chồng đang cố giữ gia đình khỏi chông chênh.}
\textit{(Each chapter in this book has three layers. The first is a story close to the life of a teacher, father, and husband trying to keep his family from instability.)}

Lớp thứ hai là một điển tích hoặc lời dạy của người xưa để soi sáng câu chuyện vừa đọc. Lớp thứ ba là bài học ngắn và câu tiếng Anh dễ nhớ để người đọc có thể dừng lại, nhẩm lại, rồi mang theo vào đời sống thường ngày.
\textit{(The second layer is a classical anecdote or teaching from the past that sheds light on the story. The third is a short lesson and an easy English line so the reader can pause, repeat it, and carry it into daily life.)}

Cuốn sách này phù hợp nhất nếu đọc chậm: mỗi ngày một mục, mỗi tối một vài trang, hoặc mỗi lúc vừa muốn mua một món không thật sự cần. Đọc ít nhưng nghĩ sâu sẽ có ích hơn đọc nhanh cho hết.
\textit{(This book works best when read slowly: one section a day, a few pages each evening, or whenever you feel tempted to buy something unnecessary. Reading little but thinking deeply is more useful than finishing quickly.)}
\end{truyen}

\begin{baihoc}
Sách tự răn chỉ có ích khi người đọc chịu dừng lại đối diện với chính thói quen của mình.
\textit{(A book of self-discipline becomes useful only when the reader is willing to face his own habits.)}
\end{baihoc}

\begin{ghinhoanh}
\textbf{Reading tips:}

Read slowly.

Pause honestly.

Bring one lesson into daily life.
\end{ghinhoanh}

\ngancach